\section{Numerical Integration and Differentiation}
\subsection{Problem 8.1}
\subsubsection*{$(a)$ Compute the approximate value of $\int_{0}^{1} x^3\:\mathrm{d}x$, first by the midpoint rule and then by the trapezoid rule.}

\textbf{Solution.} With the midpoint rule,
\begin{align*}
M &= \left(b-a\right)\:f\left(\frac{a+b}{2}\right)\\
M &= \left(1-0\right)\left(\frac{1}{2}\right)^3\\
\Aboxed{M &= \frac{1}{8}}
\end{align*}
With the trapezoid rule, we have
\begin{align*}
T &= \frac{b-a}{2}\left[f(a)+f(b)\right]\\
T &= \left(\frac{1-0}{2}\right)\left[1^3+0^3\right]\\
\Aboxed{T &= \frac{1}{2}}
\end{align*}

\subsubsection*{$(b)$ Use the difference between these two results to estimate the error in each of them.}

\textbf{Solution.} The error in computation is given by
\begin{align*}
E &\approx \frac{T-M}{3}\\
E &\approx \frac{\frac{1}{2}-\frac{1}{8}}{3}\\
\Aboxed{E &\approx \frac{1}{8}}
\end{align*}

The error in M is approximately 1/8 while error in T is approximately -1/4. This follows from the fact the the error in T is twice the error in M.

\subsubsection*{$(c)$ Combine the two results to obtain the Simson's rule approximation to the intergral.}
\textbf{Solution.} Simpson's rule is obtained by
\begin{align*}
S &= \frac{2}{3}M + \frac{1}{3}T\\
S &= \frac{2}{3}\left(\frac{1}{8}\right) + \frac{1}{3}\left(\frac{1}{2}\right) = \frac{1}{12} + \frac{1}{6}\\
\Aboxed{S &= \frac{1}{4}}
\end{align*}

\subsubsection*{$(d)$ Would you expect the latter to be exact for this problem? Why?}
\textbf{Solution.} The integral is actually
\begin{equation*}
I = \int_{0}^{1} x^3 \mathrm{d}x = \frac{x^4}{4} \Big|_0^1 = \boxed{\frac{1}{4}}
\end{equation*}
And $I=S$. The error of Simpson's method depends of the fourth and higher derivatives of the integrand that will vanish in cubic polynomials. We expect Simpson's rule to be exact in this problem.

\subsection{Problem 8.2}
\subsubsection*{$(a)$ Using the composite midpoint quadrature rule, compute the approximate value for the integral $\int_{0}^{1}\:x^3 \mathrm{d}x$, using a mesh size (subinterval length) of $h = 0.5$ and also using a mesh size of $h=1$.}
\textbf{Solution.} Using mesh size of $h=1$,
\begin{equation*}
M(h) = \left(1-0\right)\left(\frac{1}{2}\right)^3 = \frac{1}{8}  = \boxed{0.125}
\end{equation*}
Using mesh size of $h=0.5$,
\begin{align*}
M\left(\frac{h}{2}\right) = \left(0.5-0\right)\left(0.25\right)^3 + \left(1-0.5\right)\left(0.75\right)^3 = \boxed{0.21875}
\end{align*}

\subsubsection*{$(b)$ Based on the two approximate values computed  in part $(a)$, use Richardson extrapolation to compute a more  accurate approximation to the integral.}
\textbf{Solution.} In this particular example, let $p = 2$, which is the order of accuracy of the midpoint quadrature rule. The extrapolated value is given by
\begin{align*}
a_0 &= M(h) + \frac{M(h)-M(h/2)}{2^{-2}-1}\\
&= 0.125 + \frac{0.125 - 0.21875}{0.25-1}\\
&= 0.125 + \frac{-0.09375}{-0.75}\\
\Aboxed{a_0 &= 0.25}
\end{align*}

\subsubsection*{$(c)$ Would you expect the extrapolated result computed in part $(b)$ to be exact in this case? Why?}
\textbf{Solution.} In this particular extrapolation, $p=2$ and $r=4$ which is based on the error analysis of the midpoint rule. The accuracy of the extrapolation will be greater than the accuracy of the midpoint method. In this case, it is $r=4$. This really gives exact results in cubic polynomials as integrands.

\subsection{Problem 8.4}
\subsubsection*{Fill in the details of the derivation of the error estimates for the midpoint and trapezoid quadrature rules given in Section 8.3.1. In par­ticular, show that the odd-order terms drop out in both cases, as claimed.}

\textbf{Solution.} The error in the midpoint rule can be estimated using a Taylor series expansion about the midpoint $m = (b+a)/2$ of the interval $[a,b]$
\begin{align}
f(x) &= f(m) + f'(m)(x-m) + f''(m)\frac{(x-m)^2}{2} + f^{(3)}(m)\frac{(x-m)^3}{6}\\
&+ f^{(4)}(m)\frac{(x-m)^4}{24} + f^{(5)}(m)\frac{(x-m)^5}{120} + f^{(6)}(m)\frac{(x-m)^6}{720} + \cdots
\end{align}

We let $I = \int_b^a f(x)\:\mathrm{d}x$. Integrating both sides from $a$ to $b$:
\begin{align}\label{taylorint}
I(f) &= f(m)(b-a) + f'(m)\frac{(x-m)^2}{2}\Big|_a^b + f''(m)\frac{(x-m)^3}{6}\Big|_a^b + f^{(3)}(m)\frac{(x-m)^4}{24}\Big|_a^b\\
&+ f^{(4)}(m)\frac{(x-m)^5}{120}\Big|_a^b + f^{(5)}(m)\frac{(x-m)^6}{720}\Big|_a^b + f^{(6)}(m)\frac{(x-m)^7}{5040}\Big|_a^b + \cdots
\end{align}
Consider evaluating
\begin{align*}
(x-m)^n\Big|_a^b &= (b-m)^n - (a-m)^n\\
&= \left(b-\frac{a+b}{2}\right)^n - \left(a-\frac{a+b}{2}\right)^n\\
&= \left(\frac{b-a}{2}\right)^n - \left(\frac{a-b}{2}\right)^n\\
&= \frac{1}{2^n}\left[(b-a)^n-(-1)^n(b-a)^n\right]\\
\end{align*}
Looking at possible values of $n>0$, it can be shown that
\begin{equation}\label{x-m}
\boxed{
(x-m)^n\Big|_a^b = \begin{dcases}
0 &\text{if } n \text{ is even}\\
\frac{(b-a)^n}{2^{n-1}} &\text{if } n \text{ is odd}
\end{dcases}
}
\end{equation}
Applying equation (\ref{x-m}) to equation (\ref{taylorint}),
\begin{align*}
I(f) &= f(m)(b-a) + \frac{f'(m)}{2}\cdot 0 + \frac{f''(m)}{6} \frac{(x-m)^3}{2^2} + \frac{f^{(3)}(m)}{24}\cdot 0\\
&+ \frac{f^{(4)}(m)}{120} \frac{(x-m)^5}{2^4} + \frac{f^{(5)}(m)}{720}\cdot 0 + \frac{f^{(6)}(m)}{5040}\frac{(b-a)^7}{2^6} + \cdots\\
\end{align*}
Notice that the odd-order terms will drop out, yielding
\begin{align}\label{midpoint}
I(f) &= f(m)(b-a) + \frac{f''(m)}{24}(b-a)^3 + \frac{f^{(4)}(m)}{1920}(b-a)^5 + \cdots \\
\Aboxed{I(f) &= M(f) + E(f) + F(f)}
\end{align}
Where $E$ and $F$ represent the first two terms in the error expansion for the midpoint rule. To derive a comparable error expansion in the trapezoid quadrature rule, we substitute $x=a$ and $x=b$ into the Taylor series.
\begin{align*}
f(a) &= f(m) + f'(m)(a-m) + f''(m)\frac{(a-m)^2}{2} + f^{(3)}(m)\frac{(a-m)^3}{6}\\
&+ f^{(4)}(m)\frac{(a-m)^4}{24} + f^{(5)}(m)\frac{(a-m)^5}{120} + f^{(6)}(m)\frac{(a-m)^6}{720} + \cdots\\
f(b) &= f(m) + f'(m)(b-m) + f''(m)\frac{(b-m)^2}{2} + f^{(3)}(m)\frac{(b-m)^3}{6}\\
&+ f^{(4)}(m)\frac{(b-m)^4}{24} + f^{(5)}(m)\frac{(b-m)^5}{120} + f^{(6)}(m)\frac{(b-m)^6}{720} + \cdots\\
\end{align*}
Add the two series together, yielding
\begin{align}\label{taylorab}
f(a) + f(b) &= 2f(m) + f'(m)\left[(a-m)+(b-m)\right] + \frac{f''(m)}{2}\left[(a-m)^2+(b-m)^2\right] \\
&+ \frac{f^{(3)}(m)}{6}\left[(a-m)^3+(b-m)^3\right] + \frac{f^{(4)}(m)}{24}\left[(a-m)^4+(b-m)^4\right] \\
&+ \frac{f^{(5)}(m)}{120}\left[(a-m)^5+(b-m)^5\right] + \frac{f^{(6)}(m)}{720}\left[(a-m)^6+(b-m)^6\right]
\end{align}
Now, consider evaluating
\begin{align*}
(a-m)^n + (b-m)^n &= \left(a-\frac{b+a}{2}\right)^n + \left(b-\frac{b+a}{2}\right) \\
&= \left(\frac{a-b}{2}\right)^n + \left(\frac{b-a}{2}\right)^n\\
&= \frac{1}{2^n}\left[(a-b)^n+(b-a)^n\right]\\
&= \frac{1}{2^n}\left[(a-b)^n+(-1)^n(a-b)^n\right]
\end{align*}
Looking at possible values of $n>0$, it can be shown that
\begin{equation}\label{taylorabi}
\boxed{
	(a-m)^n + (b-m)^n =
	\begin{dcases}
	0 & \text{if } n \text{ is odd}\\
	\frac{(a-b)^n}{2^{n-1}} = \frac{(b-a)^n}{2^{n-1}} &\text{if } n \text{ is even}
	\end{dcases}
}
\end{equation}
Applying equation (\ref{taylorabi}) to equation (\ref{taylorab}),
\begin{align*}
f(a) + f(b) &= 2f(m) + f'(m)\cdot 0 + \frac{f''(m)}{2}\frac{(b-a)^2}{2} + \frac{f^{(3)}(m)}{6}\cdot 0\\
&+ \frac{f^{(4)}(m)}{24}\frac{(b-a)^4}{2^3} + \frac{f^{(5)}(m)}{120}\cdot 0 + \frac{f^{(6)}(m)}{720}\left[(a-m)^6+(b-m)^6\right] + \cdots \\
\end{align*}
Notice again that the odd-order terms will drop out, yielding
\begin{align*}
f(a) + f(b) &= 2f(m) + \frac{f''(m)}{4}(b-a)^2 + \frac{f^{(4)(m)}}{192}(b-a)^4 + \cdots\\
f(m) &= \frac{f(a)+f(b)}{2} - \left(\frac{f''(m)}{8}(b-a)^2 + \frac{f^{(4)}(m)}{384}(b-a)^4 + \cdots \right)\\
\Aboxed{f(m) &= T(f) - \left(\frac{f''(m)}{8}(b-a)^2 + \frac{f^{(4)}(m)}{384}(b-a)^4 + \cdots \right)}
\end{align*}
Substituting $f(m)$ in the midpoint expansion given in equation (\ref{midpoint}),
\begin{align*}
I(f) &= \left\lbrace T(f)-\left(\frac{f''(m)}{8}(b-a)^2 + \frac{f^{(4)}(m)}{384}(b-a)^4 + \cdots \right)\right\rbrace(b-a) +  \frac{f''(m)}{24}(b-a)^3 + \frac{f^{(4)}(m)}{1920}(b-a)^5 + \cdots\\
I(f) &= T(f)(b-a) + \left(\frac{-1}{8}+\frac{1}{24}\right)\cdot f''(m)(b-a)^3 + \left(\frac{-1}{384}+\frac{1}{1920}\right)\cdot f^{(4)}(m)(b-a)^5 + \cdots\\
I(f) &= T(f)(b-a) - \frac{1}{12}f''(m)(b-a)^3 - \frac{1}{480}f^{(4)}(m)(b-a)^5 + \cdots\\
\Aboxed{I(f) &= T(f) - 2E(f) - 4F(f) - \cdots}
\end{align*}
Note that
\begin{align*}
T(f) - M(f) &= T(f) - \left\lbrace I(f) - \left(\frac{f''(m)}{24}(b-a)^3 + \frac{f^{(4)}(m)}{1920}(b-a)^5+ \cdots\right)\right\rbrace\\
T(f) - M(f) &= T(f) - \left\lbrace \left(T(f)-2E(f)-4F(f)-\cdots\right) - \left(E(f) + F(f) + \cdots \right)\right\rbrace\\
\Aboxed{T(f) - M(f) &= 3E(f) + 5F(f) + \cdots}
\end{align*}
and hence the difference between the two quadrature rules provides an estimate for the dominant term in their error expansions,
\begin{equation}
\boxed{E(f) \approx \frac{T(f)-M(f)}{3}}
\end{equation}
provided that the length of the interval is sufficiently small that $(b-a)^5 \ll (b-a)^3$, and the integrand $f$ is such that $f^{(4)}$ is well-behaved. \qedsymbol

\subsection{Problem 8.6}
\subsubsection*{Suppose that Lagrange interpolation at a given set of nodes $x_1,\ldots,x_n$ is used to derive a quadrature rule. Prove that the corresponding weights are given by the integrals of the Lagrange basis functions, $w_i = \int_{a}^{b}\ell_i(x)\:\mathrm{d}x, \:i = 1,\ldots, n.$}

\textbf{Solution.} Using the method of undetermined coefiicients, we choose the weights so that the rule integrates the first $n$ polynomial basis functions exactly, which results in a system of $n$ equations in $n$ unknowns with the Lagrange basis, for example, this method results in the system of moment equations

\begin{align*}
w_1y_1\ell_1(t_1) + w_2y_1\ell_1(t_2) + \cdots + w_ny_1\ell_1(t_n) &= \int_{a}^{b}y_1\ell_1(t)\:\mathrm{d}t = y_1\int_{a}^{b}\ell_1(t)\:\mathrm{d}t \\
w_1y_2\ell_2(t_1) + w_2y_2\ell_2(t_2) + \cdots + w_ny_2\ell_2(t_n) &= \int_{a}^{b}y_2\ell_2(t)\:\mathrm{d}t = y_2\int_{a}^{b}\ell_2(t)\:\mathrm{d}t \\
&\vdots\\
w_1y_n\ell_n(t_1) + w_2y_n\ell_2(t_2) + \cdots + w_ny_n\ell_n(t_n) &= \int_{a}^{b}y_n\ell_n(t)\:\mathrm{d}t = y_n\int_{a}^{b}\ell_n(t)\:\mathrm{d}t \\
\end{align*}

Writing this system in matrix form, we have
\begin{equation*}
\begin{bmatrix}
y_1\\
y_2\\
\vdots\\
y_n
\end{bmatrix}^T
\begin{bmatrix}
\ell_1(t_1) & \ell_1(t_2) & \cdots & \ell_1(t_n) \\
\ell_2(t_1) & \ell_2(t_2) & \cdots & \ell_2(t_n) \\
\vdots & \vdots & \ddots & \vdots\\
\ell_n(t_1) & \ell_n(t_2) & \cdots & \ell_n(t_n) \\
\end{bmatrix}
\begin{bmatrix}
w_1\\
w_2\\
\vdots\\
w_n
\end{bmatrix}
=
\begin{bmatrix}
y_1\\
y_2\\
\vdots\\
y_n
\end{bmatrix}^T
\begin{bmatrix}
\int_{a}^{b}\ell_1(t)\:\mathrm{d}t\\
\int_{a}^{b}\ell_2(t)\:\mathrm{d}t\\
\vdots\\
\int_{a}^{b}\ell_n(t)\:\mathrm{d}t
\end{bmatrix}
\end{equation*}
Multiply vector $\left[\frac{1}{y_1} \frac{1}{y_2} \cdots \frac{1}{y_n}\right]$ to both sides of the equation, yielding
\begin{align*}
\begin{bmatrix}
\ell_1(t_1) & \ell_1(t_2) & \cdots & \ell_1(t_n) \\
\ell_2(t_1) & \ell_2(t_2) & \cdots & \ell_2(t_n) \\
\vdots & \vdots & \ddots & \vdots\\
\ell_n(t_1) & \ell_n(t_2) & \cdots & \ell_n(t_n) \\
\end{bmatrix}
\begin{bmatrix}
w_1\\
w_2\\
\vdots\\
w_n
\end{bmatrix}
&=
\begin{bmatrix}
\int_{a}^{b}\ell_1(t)\:\mathrm{d}t\\
\int_{a}^{b}\ell_2(t)\:\mathrm{d}t\\
\vdots\\
\int_{a}^{b}\ell_n(t)\:\mathrm{d}t
\end{bmatrix}\\
\mathbf{Lw} &= \mathbf{i}
\end{align*}
Since the property of Lagrange basis states that
\begin{equation*}
\ell_j(t_i) = \begin{cases}
1 & \text{if } i=j\\
0 & \text{if } i\neq j
\end{cases}, \:\: i=1,\ldots,n
\end{equation*}
The square matrix $\mathbf{L}$ is just an identity matrix, hence we have a solution
\begin{align*}
\mathbf{w} &= \mathbf{i}\\
\begin{bmatrix}
w_1\\
w_2\\
\vdots\\
w_n
\end{bmatrix}
&=
\begin{bmatrix}
\int_{a}^{b}\ell_1(t)\:\mathrm{d}t\\
\int_{a}^{b}\ell_2(t)\:\mathrm{d}t\\
\vdots\\
\int_{a}^{b}\ell_n(t)\:\mathrm{d}t
\end{bmatrix}\\
\Aboxed{w_i &= \int_{a}^{b}\ell_i(x)\:\mathrm{d}x, \:\:\:i = 1,\ldots, n.}
\end{align*}
Which are the corresponding weights, hence proving the claim. \qedsymbol